% ============================================================
\part{Hardwarebezug}
% ============================================================

\chapter{Wir wollen den Arduino UNO Q unter Linux finden}

\kapitelziel{
Du kannst den Arduino UNO Q unter Linux identifizieren,
den richtigen Gerätepfad finden und per udev-Regel
den stabilen Namen \texttt{/dev/arduino} einrichten.
}
\kapitelstruktur

\section{Situation}

Der Arduino UNO Q ist per USB am Raspberry Pi angeschlossen.
\texttt{sensor\_bridge.py} öffnet \texttt{/dev/arduino} --
aber die Datei existiert nicht. War es \texttt{ttyUSB0}?
Oder \texttt{ttyACM0}? Und morgen, wenn noch ein
USB-Gerät angesteckt ist, vielleicht \texttt{ttyACM1}?

\begin{aufgabeBox}
Finde den Gerätepfad des Arduino UNO Q und richte per
udev-Regel den stabilen Namen \texttt{/dev/arduino} ein.
\end{aufgabeBox}

\begin{problemBox}
Linux vergibt Gerätepfade dynamisch beim Einstecken.
\texttt{ttyACM0} kann morgen \texttt{ttyACM1} sein -- je nachdem,
was zuerst eingesteckt wird. Ohne stabile Benennung muss
jeder Entwickler den Port manuell suchen.
\end{problemBox}

\begin{haubeBox}
\textbf{Warum heißt es \texttt{/dev/ttyACM0}?}

\texttt{tty} kommt von \textit{Teletype} -- uralten Fernschreibern,
die über serielle Leitungen mit Computern verbunden waren.
Jeder serielle Port in Linux, egal ob echte Hardware oder USB,
wird als \texttt{tty}-Gerät modelliert.

\texttt{ACM} steht für \textit{Abstract Control Model} --
das USB-Klasse-Protokoll, das der Arduino-Bootloader verwendet.
Deshalb erscheint der Arduino als \texttt{ttyACM},
nicht als \texttt{ttyUSB}: er spricht ein anderes Protokoll.
\end{haubeBox}

\section{Gerät identifizieren}

\begin{lstlisting}[style=bash]
# Arduino anstecken, dann:
dmesg | tail -15
# [...] cdc_acm 1-1.4:1.0: ttyACM0: USB ACM device
# -> Port ist /dev/ttyACM0

# USB-Geraete anzeigen
lsusb
# Bus 001 Device 005: ID 2341:0069 Arduino SA
# Vendor-ID: 2341 (Arduino SA)
# Product-ID: 0069 (UNO Q / R3)

# Details fuer udev-Regel
udevadm info /dev/ttyACM0 | grep -E "ID_VENDOR_ID|ID_MODEL_ID|ID_SERIAL"
# ID_VENDOR_ID=2341
# ID_MODEL_ID=0069
# ID_SERIAL=Arduino_UNO_Q_...

# Serielle Verbindung testen (9600 Baud wie im Sketch)
stty -F /dev/ttyACM0 9600 raw
cat /dev/ttyACM0
# R=87 cm  L=124 cm  <- Arduino sendet!
# Strg+C zum Beenden
\end{lstlisting}

\section{Stabile Benennung mit udev}

\begin{lstlisting}[style=text]
# /etc/udev/rules.d/99-arduino.rules
# Arduino UNO R3 und UNO Q immer als /dev/arduino

SUBSYSTEM=="tty", \
ATTRS{idVendor}=="2341", \
ATTRS{idProduct}=="0069", \
SYMLINK+="arduino", \
MODE="0666"

# Fuer Arduino UNO R3 (ATmega328P) falls andere Product-ID:
SUBSYSTEM=="tty", \
ATTRS{idVendor}=="2341", \
ATTRS{idProduct}=="0043", \
SYMLINK+="arduino", \
MODE="0666"
\end{lstlisting}

\begin{lstlisting}[style=bash]
# Regel aktivieren
sudo nano /etc/udev/rules.d/99-arduino.rules
sudo udevadm control --reload-rules
sudo udevadm trigger

# Pruefen (Arduino muss eingesteckt sein)
ls -la /dev/arduino
# /dev/arduino -> ttyACM0

# Benutzer zur dialout-Gruppe hinzufuegen (einmalig)
sudo usermod -aG dialout $USER
# Neu einloggen!

# Jetzt funktioniert:
python3 stufe2_ros2/src/sensor_bridge.py
\end{lstlisting}

\begin{projektBox}
\textbf{Bisher:} \texttt{sensor\_bridge.py} öffnet
\texttt{/dev/arduino} -- der Pfad existierte nicht.\\
\textbf{Heute:} udev-Regel eingerichtet.
\texttt{/dev/arduino} zeigt immer auf den richtigen Port,
egal in welcher Reihenfolge Geräte eingesteckt werden.\\
\textbf{Auch in docker-compose.yml:}
\texttt{devices: ["/dev/arduino:/dev/arduino"]} funktioniert jetzt.
\end{projektBox}

\begin{merksatzBox}
udev identifiziert Geräte über USB-IDs, nicht über den
zufälligen Gerätepfad. Vendor-ID \texttt{2341} gehört
immer zu Arduino SA -- egal auf welchem Raspberry Pi
und egal welcher Port vergeben wird.
\end{merksatzBox}

\begin{fehlerBox}
\begin{itemize}
  \item \textbf{Permission denied:}
        \texttt{sudo usermod -aG dialout \$USER} -- dann neu einloggen.
  \item \textbf{/dev/arduino existiert nicht nach Neustart:}
        Regel-Datei hat Tippfehler. Prüfen mit
        \texttt{udevadm test \$(udevadm info -q path /dev/ttyACM0)}.
\end{itemize}
\end{fehlerBox}

\begin{robotikBox}
ROS2-Treiber für Lidar-Scanner, Kameras und IMUs installieren
automatisch udev-Regeln. Das \texttt{ros-jazzy-rplidar-ros}-Paket
legt \texttt{/dev/rplidar} an -- nach demselben Prinzip.
Wer es einmal selbst eingerichtet hat, versteht warum
ROS2-Hardware \enquote{einfach funktioniert}.
\end{robotikBox}

% ── Kapitel 24 ────────────────────────────────────────────────
\chapter{Wir wollen den Differentialantrieb von Python aus steuern}

\kapitelziel{
Du verstehst das Kommunikationsprotokoll zwischen
differentialantrieb.ino und Python und implementierst
die Motorsteuerung für \texttt{stufe2\_ros2/}.
}
\kapitelstruktur

\section{Situation}

\texttt{differentialantrieb.ino} fährt zurzeit blind
geradeaus. Die Hinderniserkennung ist getrennt.
In Stufe~2 soll Python die Intelligenz übernehmen:
Sensordaten empfangen, entscheiden, Motorbefehle senden.

\begin{aufgabeBox}
Erweitere \texttt{differentialantrieb.ino} um einen
Serial-Befehlsempfänger und implementiere die passende
Python-Klasse \texttt{MotorBridge} in \texttt{stufe2\_ros2/}.
\end{aufgabeBox}

\begin{wissenBox}
Wir brauchen: das bestehende L298N-Protokoll aus
\texttt{differentialantrieb.ino}, Serial-Kommunikation
und das Prinzip: Python = Intelligenz, Arduino = Echtzeit.
\end{wissenBox}

\section{Das bestehende L298N-Interface}

Aus \texttt{differentialantrieb.ino} kennen wir bereits:

\begin{lstlisting}[style=text]
Motor A (links):  ENA=D5(PWM), IN1=D4, IN2=D7
Motor B (rechts): ENB=D6(PWM), IN3=D8, IN4=D12

Pinbelegung aus docs/Verdrahtung.md:
  vorwaerts(n):   IN1=H IN2=L IN3=H IN4=L, ENA=ENB=n
  rueckwaerts(n): IN1=L IN2=H IN3=L IN4=H, ENA=ENB=n
  links(n):       IN1=L IN2=H IN3=H IN4=L, ENA=ENB=n
  rechts(n):      IN1=H IN2=L IN3=L IN4=H, ENA=ENB=n
  stopp():        alle LOW, ENA=ENB=0
\end{lstlisting}

\section{Arduino-Erweiterung: Befehle empfangen}

\begin{lstlisting}[style=text]
// Ergaenzung in differentialantrieb.ino:
// Am Ende von loop() hinzufuegen:

void loop() {
  if (kontaktAusgeloest) {
    ausfuehrenAusweichmanoever();
    return;
  }

  // NEU: Serial-Befehle von Python empfangen
  if (Serial.available() > 0) {
    String cmd = Serial.readStringUntil('\n');
    cmd.trim();

    if (cmd == "S")          { stopp(); }
    else if (cmd == "V")     { vorwaerts(SPEED_NORMAL); }
    else if (cmd == "R")     { rueckwaerts(SPEED_RUECKWAERTS); }
    else if (cmd == "L")     { links(SPEED_DREHEN); }
    else if (cmd == "D")     { rechts(SPEED_DREHEN); }
    else if (cmd.startsWith("P ")) {
      // "P links rechts" -- individuelle PWM-Werte
      int sp = cmd.indexOf(' ');
      int lp = cmd.indexOf(' ', sp+1);
      int l  = cmd.substring(sp+1, lp).toInt();
      int r  = cmd.substring(lp+1).toInt();
      // Motorsteuerung mit individuellen Werten...
      vorwaerts(l); // vereinfacht
    }
    Serial.println("OK");
  }
  else {
    // Kein Befehl: Stopp (Python entscheidet)
    stopp();
  }
}
\end{lstlisting}

\section{Python-Seite: MotorBridge}

\begin{lstlisting}[style=python]
# stufe2_ros2/src/motor_bridge.py
import serial, time

class MotorBridge:
    BEFEHLE = {
        'stop':      'S',
        'vor':       'V',
        'zurueck':   'R',
        'links':     'L',
        'rechts':    'D',
    }

    def __init__(self, port='/dev/arduino', baud=9600):
        self.ser = serial.Serial(port, baud, timeout=1)
        time.sleep(2)  # Arduino Neustart abwarten
        print(f"[motor_bridge] Verbunden: {port}")

    def sende(self, befehl):
        cmd = self.BEFEHLE.get(befehl, 'S')
        self.ser.write(f"{cmd}\n".encode())
        antwort = self.ser.readline().decode().strip()
        return antwort == "OK"

    def navigiere(self, rechts_cm, links_cm):
        minimum = min(rechts_cm, links_cm)
        if minimum < 25:
            self.sende('stop')
        elif minimum < 50:
            self.sende('vor')   # wird langsamer in Band 2
        elif rechts_cm < links_cm:
            self.sende('links')
        else:
            self.sende('vor')

    def __del__(self):
        if hasattr(self, 'ser') and self.ser.is_open:
            self.sende('stop')
            self.ser.close()

# Verwendung:
# motor = MotorBridge('/dev/arduino')
# motor.navigiere(rechts_cm=87, links_cm=124)
\end{lstlisting}

\begin{projektBox}
\textbf{Heute:} \texttt{stufe2\_ros2/src/motor\_bridge.py} --
Python steuert die Motoren über das Extended Serial-Protokoll.\\
\textbf{Die Kette ist vollständig:}
\texttt{hinderniserkennung.ino} $\to$ \texttt{sensor\_bridge.py}
$\to$ \texttt{SensorBus} $\to$ \texttt{motor\_bridge.py}
$\to$ \texttt{differentialantrieb.ino (erweitert)} $\to$ L298N.\\
\textbf{Nächstes Kapitel:} Warum dieser Stack jetzt ROS2 braucht.
\end{projektBox}

\begin{merksatzBox}
Python = Intelligenz (Entscheidungslogik, Konfiguration, Logging).
Arduino = Echtzeit (PWM, Interrupts, Motoransteuerung).
Diese Trennung ist kein Umweg -- sie ist Systemarchitektur.
Jede Schicht macht, was sie am besten kann.
\end{merksatzBox}

\begin{robotikBox}
In Band~2 ersetzt \texttt{motor\_bridge.py} den direkten
Serial-Aufruf durch:
\begin{lstlisting}[style=python]
cmd = Twist()
cmd.linear.x = 0.3   # m/s vorwaerts
self.cmd_vel_pub.publish(cmd)
\end{lstlisting}
Ein Arduino-seitiger ROS2-Serial-Bridge-Node
(z.\,B. \texttt{rosserial} oder \texttt{micro\_ROS})
übersetzt den Twist-Befehl in PWM.
Das Prinzip -- Python schickt Befehle, Arduino führt aus --
bleibt exakt dasselbe.
\end{robotikBox}
